\documentclass[../main.tex]{subfiles}
\begin{document}
%==========================================TIÊU ĐỀ TẬP========================================
\begin{table}[h]
    \begin{adjustwidth}{-1cm}{}
        \begin{tabular}{>{\centering\arraybackslash}p{6cm}|>{\raggedright\arraybackslash}p{12.5cm}}
            \multirow{5}{6cm}
            % Cột bên trái: ÉPISODE và số tập
            {\\[-40pt]
                \flaregothic\fontsize{20pt}{20pt}\selectfont \centering
                \color{eptype}HISTOIRE\color{black}\\
                \freshbot\fontsize{64pt}{64pt}\selectfont
                \color{epnum}B\\[6pt]
                \cafeteria\fontsize{20pt}{20pt}\selectfont\color{epnum}(D-B)\color{black}
                \\[18pt]
                % \sen\fontsize{18pt}{18pt}\selectfont
                % \color{epattr}Interstitiel
                \color{black}
            }
             &
            % Dòng thứ nhất: tên tập tiếng Nhật
            {\fotskip\fontsize{18pt}{18pt}\selectfont \ruby{小}{ちい}さな\ruby{レ}{れ}\ruby{ン}{ん}\ruby{ズ}{ず}の\ruby{向}{む}こう\ruby{側}{gaw}
            }         \\[6pt]
            % Dòng thứ hai: tên tập tiếng Anh
             & {\cafeteria\fontsize{18pt}{18pt}\selectfont The World Through a Small Lens}                                              \\[4pt]
            % Dòng thứ ba: tên tập tiếng Việt
             & {\vnmsans\fontsize{18pt}{18pt}\selectfont Thế giới qua tấm kính nhỏ}                                                     \\[8pt]
            % Dòng thứ tư: Nhân vật xuất hiện
             & {\caviar{\fontsize{14pt}{14pt}\selectfont Nhân vật: \emp{kuon2} \emp{ryuuhou2} \emp{kaiho2} \emp{suzuno2} \emp{naomi2}}}
            \\[8pt]
        \end{tabular}
    \end{adjustwidth}
\end{table}
%=========================================TRUYỆN CHÍNH========================================
\color{black}

\txtbx[2]{
    {\color{vol} \uline{Tóm tắt:}} Buổi kiểm tra sức khỏe định kỳ đầu năm ở trường trung học bất ngờ giáng một đòn choáng váng xuống gia đình Hikawa: trong khi mắt phải của Ryuuhou đạt thị lực tuyệt đối mười phần mười, thì mắt trái của cô bé lại tụt dốc thê thảm chỉ còn vỏn vẹn một phần mười! Thói quen nằm nghiêng trùm chăn bấm điện thoại thâu đêm suốt sáng đến hai giờ sáng rốt cuộc đã để lại hậu quả nghiêm trọng. Đứng trước sự thật bàng hoàng ấy, chị cả Kaiho nổi trận lôi đình với phản ứng vô cùng quyết liệt, trong khi bố Ryujin dù giữ thái độ hòa nhã, trầm ấm nhưng lại tỏ rõ sự nghiêm túc chưa từng thấy. Một phán quyết không thể chối từ được đưa ra: Ryuuhou phải đi khám chuyên sâu và cắt kính bằng mọi giá.
}

\stpt{-Histoire 1-}
\stcap{Con số một phần mười}

\pov{Hikawa Kuon}

Buổi chiều muộn của một ngày đầu tháng Ba, nắng xuân ở thủ đô Kawachi còn vương chút se lạnh đặc trưng của đợt gió mùa đông bắc muộn.

Tôi vừa kết thúc ca họp trực tuyến kéo dài ba tiếng đồng hồ với bộ phận kỹ thuật bên Tokyo, xoa xoa đôi mắt mỏi nhừ rồi bước từ tầng năm xuống sảnh sinh hoạt chung của tầng ba để pha một tách trà nóng. Lúc này, cả nhà đã bắt đầu lục tục trở về sau một ngày học tập và làm việc.

Căn phòng tầng ba lúc này ngập tràn ánh hoàng hôn vàng cam ấm áp chiếu qua khung cửa sổ lớn nhìn ra dải công viên ven sông. Trên chiếc bàn gỗ dài, bé út Naomi đang ngoan ngoãn ngồi vẽ tranh cùng chú chó già Mi nằm sưởi ấm dưới chân, trong khi Suzuno vừa đi chợ về, đang cẩn thận phân loại từng bó rau muống xanh mướt cho bữa cơm tối.

Đúng lúc ấy, cánh cửa gỗ bật mở. Ryuuhou bước vào phòng với dáng vẻ uể oải quen thuộc. Cô nhóc ném toẹt chiếc cặp sách nặng trịch lên góc bàn trà, tiện tay lôi từ trong tủ lạnh ra một que kem đậu xanh mát lạnh rồi ngồi phịch xuống chiếc ghế đệm bọc vải nỉ, chân gác chữ ngũ lên bàn, ngửa cổ cắn một miếng kem giòn rụm.

``Học hành gì mà trông như người mất hồn thế em?'' tôi mỉm cười, vừa rót nước sôi vào ấm trà vừa hỏi với sang.

``Mệt chết đi được ấy Onii ơi,'' Ryuuhou làu bàu qua kẽ răng, que kem vẫn ngậm trong miệng. ``Hôm nay toàn trường em bị lùa đi khám sức khỏe định kỳ. Đo chiều cao, cân nặng, rồi lại còn thử máu với đo thị lực. Đứng xếp hàng rồng rắn cả buổi chiều dưới sân thể chất, chân cẳng mỏi rã rời.''

Nói rồi, cô nàng thò tay vào túi áo khoác đồng phục, rút ra một tờ giấy khổ A4 gấp làm tư đã hơi nhàu nhĩ, tiện tay quẳng bừa lên mặt bàn trà cạnh đống bút màu của Naomi: ``Phiếu kết luận sức khỏe của em đấy. Bác sĩ bảo chiều cao tăng được có nửa phân, chán òm.''

Tôi chưa kịp bước lại gần cầm tờ giấy lên xem thì cánh cửa hành lang lại lách cách mở ra lần nữa. Kaiho vừa tan ca ở xưởng cơ khí trở về. Cô em gái song sinh của tôi vẫn còn mặc chiếc áo khoác bảo hộ mỏng màu xanh công nhân, trên ngón tay vương chút vệt dầu máy chưa kịp rửa sạch. Vừa bước vào phòng, tính cách ngăn nắp và nghiêm cẩn của Kaiho lập tức phát huy tác dụng khi cô nhìn thấy tờ giấy nằm vắt vẻo bên mép bàn trà: ``Về nhà là vứt đồ đạc lung tung thế hả Ryuu? Giấy tờ trường phát mà cứ như giấy lộn thế này.''

Kaiho vừa càu nhàu vừa cúi người nhặt tờ phiếu khám sức khỏe lên, định bụng gấp lại gọn gàng để kẹp vào sổ theo dõi gia đình.

Thế nhưng, chỉ vừa liếc mắt lướt qua những dòng kết luận y khoa in mực đen, đôi con ngươi màu xanh của cô em song sinh bỗng nhiên co thắt lại dữ dội. Bước chân của em ấy khựng lại giữa sàn nhà. Bàn tay đang cầm tờ giấy siết chặt đến mức mép giấy hơi nhăn nhúm.

``\dots Cái gì thế này?!''

Giọng của Kaiho không hề lớn, nhưng sắc lạnh và gãy gọn đến mức làm cho cả căn phòng bỗng chốc rơi vào một khoảng lặng kỳ lạ. Suzuno đang nhặt rau ở góc bếp cũng giật mình ngẩng đầu lên, còn Naomi thì tròn xoe mắt ngơ ngác nhìn sang chị cả.

Tôi nhướng mày, cảm thấy có điều gì đó chẳng lành: ``Sao thế Kaiho? Huyết áp hay tim mạch của con bé có vấn đề à?''

``Anh lại đây mà nhìn này!'' Kaiho quay ngoắt người lại, chìa thẳng tờ giấy vào mặt tôi, ngón tay trỏ run run gõ mạnh vào khung kiểm tra chuyên khoa mắt.

Tôi đón lấy tờ phiếu, mắt dán vào hai dòng chữ in hoa rõ mồn một bên dưới mục \textit{Thị lực không kính}:

\begin{itemize}
    \item \textbf{Mắt phải (V.D):} \textit{10/10}
    \item \textbf{Mắt trái (V.S):} \textit{1/10}
\end{itemize}

Tôi đứng sững sờ. Não bộ mất đúng ba giây để xử lý hai con số đối nghịch đến hoang đường ấy.

Mười phần mười\dots\ và một phần mười?!

Tôi bất giác ngẩng lên, nhìn thẳng vào con bé với vẻ mặt dở khóc dở cười, buông một lời bình thản nhưng phũ phàng đến tận cùng:

``Mắt một trên mười thì coi như mù luôn rồi còn gì\dots''

``Onii-chan! Anh ăn nói kiểu gì thế hả?!'' Ryuuhou giật nảy mình, que kem suýt rơi khỏi tay, mặt đỏ tía tai cãi lại.

Nhưng con bé chưa kịp thanh minh với tôi thì cơn thịnh nộ thực sự đã giáng xuống.

Trong nhà này, nếu tôi là người điềm tĩnh quan sát thì Kaiho lại là người bảo bọc các em một cách quyết liệt nhất. Vừa nghe thấy câu nói của tôi, ngọn lửa giận trong mắt Kaiho bùng lên dữ dội. Cô nàng bước phắt tới, đập mạnh bàn tay xuống mặt bàn gỗ đánh *RẦM* một cái khiến Naomi giật bắn người, hai mắt trừng trừng như muốn bốc hỏa:

``Anh Kuon nói không sai một chữ nào đâu! Em còn dám già mồm cãi lại à Ryuu?! Bảng thị lực quốc tế người ta đo từ hàng chữ to nhất đến hàng chữ nhỏ nhất! Một phần mười tức là ngay cả cái chữ cái to tổ bố bằng bàn tay ở hàng đầu tiên em cũng không nhìn rõ đấy có biết không?! Một mắt mười, một mắt một, em tính biến mình thành người khuyết tật một bên mắt đấy hả?!''

Ryuuhou bị tiếng đập bàn của chị cả dọa cho co rúm người lại, bờ vai rụt hẳn xuống, ấp úng phân bua: ``Thì\dots\ thì bác sĩ ở trường đo vội quá nên em đọc nhầm thôi mà\dots\ Với lại mắt phải em vẫn mười phần mười, ban ngày đi đứng chạy nhảy có vấp vào ai đâu, hai mắt nó tự bù trừ cho nhau chứ bộ\dots''

``Bù trừ cái nỗi gì!'' Kaiho giận run cả người, chống hai tay bên sườn, bước tới sát chiếc ghế đệm như thể chuẩn bị tóm gáy lôi cô em gái đi trừng phạt. ``Mắt người chứ có phải bóng đèn hỏng một sợi đốt vẫn sáng le lói đâu mà đòi bù trừ!''

Đúng lúc không khí trong phòng căng thẳng như một kho thuốc súng sắp nổ tung, tiếng lách cách mở cửa từ cầu thang vang lên. Bố Ryujin vừa từ xưởng trở về, trên tay xách theo chiếc hộp đồ nghề kim loại nặng trịch.

Vừa bước vào ngưỡng cửa tầng ba, nhìn thấy Kaiho đang chống nạnh đằng đằng sát khí còn Ryuuhou thì co ro trên ghế, bố khẽ nhướng mày, cất giọng trầm ấm, hòa nhã quen thuộc:

``Có chuyện gì mà hai chị em mới chiều hôm đã ầm ĩ cả lên thế? Kaiho, con gái con lứa sao lại đập bàn đập ghế chan chát thế kia?''

``Bố! Bố vào mà xem con gái cưng của bố này!'' Kaiho hậm hực giật lấy tờ giấy khám sức khỏe trên bàn trà, bước nhanh tới đưa cho bố, giọng vẫn chưa hết gay gắt. ``Bố xem thị lực của nó đi! Mắt phải mười phần mười, mắt trái còn đúng một phần mười! Bác sĩ đóng dấu đỏ chót đây này!''

Khác với phản ứng bùng nổ quyết liệt của Kaiho, bố Ryujin không hề lớn tiếng hay nổi giận. Ông từ tốn đặt hộp đồ nghề xuống góc tường, cởi đôi găng tay lao động rồi rút chiếc kính lão giắt ở túi áo ngực đeo lên. Bố cầm lấy tờ giấy, nheo mắt đọc kỹ từng thông số.

Thế nhưng, khi ánh mắt ông dừng lại ở dòng chữ \textit{V.S: 1/10}, nụ cười hiền hậu, hòa nhã thường ngày trên gương mặt dạn dày gió sương của bố dần thu lại. Thay vào đó là một vẻ mặt đăm chiêu, nghiêm túc đến lạ thường.

Bố tháo chiếc kính lão cầm trên tay, bước lại gần bàn trà rồi kéo một chiếc ghế gỗ, ngồi xuống ngay trước mặt Ryuuhou. Ánh mắt bố nhìn thẳng vào con gái—không hề có sự trách phạt hay nạt nộ như Kaiho, nhưng sự nghiêm cẩn và từng trải trong đôi mắt người cha lại mang một sức nặng ghê gớm.

``Ryuuhou này,'' giọng bố trầm xuống, ôn tồn nhưng từng chữ đều rành rọt, dứt khoát. ``Bố làm thợ cơ khí cả đời, con biết rồi đấy. Một cỗ máy đang vận hành trơn tru, chỉ cần một con ốc lỏng hay một trục quay lệch một ly thôi là cả dây chuyền phải dừng lại, vứt đi cả tấn phôi thép. Đôi mắt con người sinh ra có đôi có cặp là để cân bằng thị giác, giúp con định lượng khoảng cách và chiều sâu của không gian. Một bên mười, một bên một—tức là mắt trái của con gần như đã ngừng hoạt động hoàn toàn. Con bảo hai mắt tự bù trừ, thực chất là con đang ép mắt phải gánh toàn bộ tải trọng, còn mắt trái thì đang dần bị đào thải. Con nghĩ chuyện này là chuyện đùa được sao?''

Lời giải thích từ tốn, hòa nhã nhưng phân tích cặn kẽ và đầy vẻ nghiêm túc của bố khiến sự bướng bỉnh cuối cùng của Ryuuhou bỗng chốc xẹp lép. Cô nhóc cúi gằm mặt xuống gối, không dám đối diện với ánh nhìn của bố.

Căn phòng tầng ba rơi vào một khoảng lặng nặng nề.

Đúng lúc ấy, Suzuno rụt rè bước lại gần bàn trà, tay vẫn còn cầm chiếc khăn lau bếp. Giọng cô bé nhỏ nhẹ nhưng từng lời nói ra lại như những nhát búa đóng đinh vào sự thật: ``Dạo gần đây\dots\ em cũng thấy mắt của Ryuuhou có vấn đề thật đấy ạ. Cuối tuần trước hai chị em đánh cầu lông ở sân sau, những quả cầu bổng em phát sang phía bên tay trái của em, em ấy toàn vung vợt hụt liên tục. Lúc đầu em cứ tưởng em ấy nhường, nhưng sau đó nhìn kỹ thì thấy em ấy toàn phải nghiêng hẳn cả đầu sang một bên mới đón được cầu\dots''

Bé út Naomi đang ôm chú chó Mi dưới sàn cũng giơ bàn tay nhỏ xíu lên, ngây thơ góp thêm một đòn chí mạng: ``Đúng rồi đó ạ! Buổi tối trùm chăn ngủ ở tầng ba, Naomi toàn thấy chị Ryuuhou trùm kín chăn từ đầu đến chân rồi bật điện thoại sáng trưng lên xem tới khuya lắc khuya lơ! Chị ấy toàn nằm nghiêng sang bên phải, áp một bên mặt xuống gối rồi nhắm tịt một mắt lại, chỉ mở đúng con mắt trái nhìn vào màn hình thôi! Giống hệt mấy chú cướp biển một mắt trong phim hoạt hình ấy ạ!''

Lời tố giác hồn nhiên của Naomi khiến mặt Ryuuhou từ trắng bệch thoắt cái chuyển sang đỏ bừng như quả cà chua chín. Con bé lúng túng giấu que kem đã chảy nước ra sau lưng, hai chân co rúm lại trên ghế đệm, hoàn toàn không thể cãi lại nửa lời trước những chứng cứ mười mươi từ chính các em gái cùng phòng.

Kaiho đứng bên cạnh nghe xong thì tức đến mức nghiến răng ken két, muốn xông lên mắng tiếp, nhưng bố Ryujin khẽ giơ bàn tay thô ráp lên, ra hiệu cho con gái cả dừng lại.

Bố nhìn thẳng vào Ryuuhou, ánh mắt chan chứa sự xót xa của người làm cha nhưng thái độ thì kiên quyết không lay chuyển:

``Nghe các em nói chưa con? Trùm chăn xem điện thoại trong bóng tối đến hai giờ sáng, lại còn đè ép một bên mắt suốt hàng tháng trời. Mắt con người bằng da bằng thịt chứ có phải bằng thép tôi nguội đâu mà chịu nổi sự đày đọa ấy? Con có biết nhược thị nguy hiểm thế nào không? Để lâu thêm chút nữa, đường dẫn truyền thị giác từ mắt trái lên não sẽ bị teo lại vĩnh viễn, lúc đó có muốn chữa cũng không kịp nữa.''

Tôi cũng ngồi xổm xuống trước mặt em gái, khẽ hỏi: ``Dạo này có hay bị nhức thái dương với nhòe chữ trên bảng không Ryuu?''

Ryuuhou cắn chặt môi, mắt rơm rớm, lí nhí thú nhận: ``Thì\dots\ thỉnh thoảng nhìn chữ trên bảng em thấy nó cứ bị nhòe đôi ra, nhức buốt nửa đầu bên trái. Em tưởng là do thiếu ngủ thôi\dots''

Bố Ryujin thở dài một hơi đầy trăn trở. Bố vỗ nhẹ lên vai Ryuuhou, cất giọng hòa nhã nhưng dứt khoát đưa ra phán quyết tối cao:

``Biết đau, biết nhức tức là cơ thể đang kêu cứu đấy con ạ. Không chần chừ gì nữa. Ngay sáng mai, Kuon với Kaiho thu xếp công việc đưa em nó đến Bệnh viện Mắt trung tâm thành phố khám chuyên sâu. Phải đo khúc xạ chuẩn rồi đi cắt kính cho con bé đàng hoàng.''

Nghe đến hai chữ ``cắt kính'', Ryuuhou lập tức giãy nảy lên như đỉa phải vôi, hai tay ôm chặt lấy đầu, mếu máo kêu la: ``Ôi không không không! Bố ơi con xin bố đấy! Đừng bắt con đeo kính mà! Một đứa chơi nhạc rock như con mà phải vác hai cái đít chai dày cộp lên mặt thì trông ngố tàu chết đi được! Lũ bạn trong ban nhạc nó cười cho thối mũi mất thôi! Con không đeo đâu, con thà mù một mắt chứ không đeo kính đâu!!''

``Ăn nói bậy bạ!'' Kaiho lại gắt lên.

Bố Ryujin khẽ lắc đầu, nụ cười hiền từ thoáng hiện lại trên khóe môi trước vẻ mè nheo trẻ con của cô con gái ngổ ngáo, nhưng giọng điệu của bố vẫn giữ nguyên sự nghiêm túc tuyệt đối:

``Xấu xí thì có khối cách làm cho bớt xấu, chứ hỏng mắt rồi thì không ai đền lại ánh sáng cho con được đâu. Bố đã quyết rồi thì không có nhưng nhị gì hết. Tối nay đưa điện thoại đây bố giữ giùm, mười giờ là cả phòng tầng ba phải tắt đèn đi ngủ. Sáng mai chín giờ Kuon và Kaiho áp giải con bé đi.''

Nói xong, bố vỗ vai tôi một cái đầy tin cậy rồi thong thả bước lên cầu thang tầng năm cất đồ nghề và chuẩn bị tắm rửa, để lại Ryuuhou ngồi bệt trên ghế đệm với vẻ mặt như thể vừa bị kết án lưu đày biệt xứ.

\stpt{-Histoire 2-}
\stcap{Liều thuốc liệt điều tiết}

Sáng thứ Bảy hôm sau, cả khu phố Anwa thức giấc trong tiết trời trong veo mát rượi.

Theo thông lệ, đây là ngày duy nhất trong tuần mà tất cả các thành viên trong gia đình Hikawa đều không phải đến công sở hay trường học. Thông thường, vào buổi sáng cuối tuần thảnh thơi như thế này, Kaiho sẽ dậy từ sớm, khoác chiếc tạp dề vải bạt màu xanh đậm rồi cùng bố Ryujin xuống xưởng cơ khí ở tầng một để tiện phôi, cân chỉnh lại các trục bánh răng và lau chùi dầu mỡ cho dàn máy cắt gọt kim loại.

Thế nhưng sáng nay, kế hoạch xưởng máy đành phải tạm gác lại. Với trọng trách ``áp giải phạm nhân'' do đích thân bố Ryujin giao phó, Kaiho đã chuẩn bị xong xuôi từ tám giờ sáng. Cô nàng diện một bộ trang phục thường nhật gọn gàng gồm áo thun đen, quần jeans lửng và giày thể thao, đứng khoanh tay bên cửa ra vào tầng ba với vẻ mặt nghiêm nghị như một viên sĩ quan giám thị.

Điều bất ngờ là cuộc hành quân sáng nay không chỉ có ba người chúng tôi.

Suzuno đã cẩn thận chuẩn bị sẵn một túi xách vải, bên trong đựng mấy chai nước lọc và khăn giấy ướt. Cô bé nhỏ nhẹ giải thích lý do đi cùng: ``Em đi cùng để đỡ Ryuuhou ạ. Bác sĩ mà nhỏ thuốc hay đo mắt xong, lỡ mắt em ấy bị mờ không nhìn rõ đường sá thì có em dắt tay cho an toàn. Với lại tiện đường ghé qua hiệu sách lớn cạnh bệnh viện, em muốn mua thêm ít tập giấy vẽ cho bé Naomi và vài cuộn chỉ thêu.''

Về phần bé út Naomi, con bé đã xỏ sẵn đôi giày búp bê màu hồng, tay ôm chặt lấy cánh tay tôi, hai mắt lấp la lấp lánh như chuẩn bị đi trẩy hội: ``Naomi cũng muốn đi cùng anhvà các chị! Ở nhà một mình với chú cún Mi buồn lắm ạ, Naomi muốn đi xem người ta đo mắt cơ!''

Thế là từ một chuyến đi khám bệnh đơn thuần, gia đình Hikawa lại biến nó thành một cuộc diễu hành rầm rộ của cả năm anh chị em. Ryuuhou bị Kaiho kẹp chặt một bên cánh tay, lê từng bước chân nặng như chì qua con dốc trước cửa nhà, vượt qua dải công viên ven sông rồi hướng thẳng về phía Bệnh viện Mắt và Kính thuốc trung tâm thành phố Kawachi.

\ct

Phòng khám mắt buổi sáng thứ Bảy tương đối vắng vẻ và yên tĩnh, thoang thoảng mùi cồn sát khuẩn nhẹ hòa lẫn hương sáp lau sàn sạch sẽ.

Sau khi làm thủ tục đăng ký, Ryuuhou được vị bác sĩ nhãn khoa trung niên dẫn vào phòng khúc xạ chuyên sâu. Căn phòng được bố trí ánh sáng tiêu chuẩn, trang bị đầy đủ các dàn máy móc hiện đại bậc nhất. Bác sĩ yêu cầu Ryuuhou đặt cằm lên bệ đỡ của chiếc máy đo khúc xạ tự động, mắt nhìn thẳng vào một ống kính tối màu.

``Nào, cháu nhìn thẳng vào bức tranh phía trước nhé. Có thấy hình một quả khinh khí cầu màu đỏ đang bay trên con đường dài hun hút không? Chớp mắt tự nhiên, đừng nheo mắt nhé,'' giọng bác sĩ trầm ấm vang lên.

*Tách! Tách! Tách!*

Những tia sáng đỏ quét nhẹ qua đồng tử của Ryuuhou. Sau vài thao tác bấm nút rành rọt, bác sĩ chuyển cô bé sang ghế đối diện với bảng thị lực điện tử chiếu trên tường phản quang ở cự ly năm mét. Một chiếc gọng kính thử nghiệm bằng kim loại nặng trịch với hàng chục khay cài thấu kính được đeo lên sống mũi con bé.

Bác sĩ dùng một tấm bìa đen che mắt phải của Ryuuhou lại, bắt đầu dùng điều khiển từ xa chỉ vào từng hàng chữ cái bên mắt trái:

``Chữ gì đây cháu?''

``Chữ\dots\ chữ C ạ.''

``Hàng dưới này thì sao?''

``Chữ E\dots\ chữ F\dots\ và chữ D ạ.''

Ryuuhou đọc chậm rãi, thi thoảng hơi nghiêng đầu nhưng tuyệt nhiên không hề có vẻ chật vật, mờ tịt đến mức không đọc nổi chữ cái to nhất như trong tờ phiếu khám sức khỏe ở trường.

Bác sĩ tiếp tục trượt xuống hàng chữ thứ bảy: ``Hàng này cháu nhìn rõ không?''

Ryuuhou chớp chớp mắt hai cái, ngập ngừng một thoáng rồi đọc chuẩn xác: ``Chữ P, chữ T, chữ O, chữ Z.''

Đến hàng thứ tám thì con bé bắt đầu lắc đầu: ``Hàng này hơi nhòe rồi ạ, cháu chỉ đoán lờ mờ được chữ đầu tiên thôi.''

Bác sĩ tháo chiếc gọng thử nghiệm xuống, ghi chép thoăn thoắt vào sổ bệnh án rồi quay ghế lại nhìn chúng tôi: ``Kết quả kiểm tra sơ bộ đây. Mắt phải của cháu đạt mười phần mười tuyệt đối, thị lực cực tốt. Còn mắt trái thì đạt khoảng bảy phần mười, hơi có dấu hiệu cận thị nhẹ chừng không phẩy bảy lăm đi-ốp.''

Tôi và Kaiho đứng tựa vào vách tường nãy giờ, vừa nghe thấy con số ``bảy phần mười'' thì đồng loạt thở phào một hơi dài như vừa trút được cả tảng đá ngàn cân đè nặng trong lồng ngực. Kaiho ôm ngực, hai vai chùng hẳn xuống, nét mặt căng thẳng từ hôm qua đến giờ mới chịu dãn ra một chút.

``Bác sĩ ơi\dots\ thế tại sao hôm qua ở trường, phiếu khám lại ghi mắt trái con bé chỉ còn có một phần mười ạ?'' tôi vội bước lên một bước, thắc mắc hỏi. ``Con số chênh lệch nhau lớn quá làm cả nhà cháu lo sốt vó suốt từ tối qua đến giờ.''

Bác sĩ bật cười hiền hậu, tháo kính xuống lau nhẹ: ``Chuyện này xảy ra thường xuyên ở lứa tuổi học sinh các cháu lắm. Hôm qua cháu thức khuya đúng không? Lại còn nằm nghiêng dùng điện thoại trong phòng tối nữa?''

Ryuuhou đỏ ửng mặt, gật đầu đánh rụt một cái.

``Đó chính là nguyên nhân đấy,'' bác sĩ giải thích cặn kẽ. ``Khi các cháu dán mắt vào màn hình cự ly quá gần trong bóng tối đặc quánh, các cơ thể mi xung quanh thể thủy tinh phải gồng lên co thắt liên tục suốt nhiều giờ liền. Khi đến trường khám, mắt chưa kịp hồi phục lại gặp môi trường đo vội vã, thiếu sáng và tâm lý căng thẳng, cơ thể mi bị co quắp cực độ dẫn đến hiện tượng `cận thị giả' cục bộ. Thị lực tức thời vì thế mà tụt dốc không phanh xuống một phần mười. Hôm nay được ngủ một giấc, cơ mi thả lỏng bớt nên mắt đã hồi lại bảy phần mười rồi.''

Nghe đến đây, Ryuuhou lập tức ngẩng phắt đầu lên, toe toét cười đắc chí nhìn Kaiho: ``Đấy! Em đã bảo là do bác sĩ ở trường đo nhầm mà chị Kaiho cứ làm ầm cả nhà lên! Mắt em bảy phần mười chứ có phải một phần mười đâu!''

``Cấm có cãi!'' Kaiho lập tức trừng mắt, cốc nhẹ một cái vào trán cô em gái. ``Bảy phần mười thì vẫn là cận thị! Vẫn là lệch thị lực so với mắt phải! Em nghe bác sĩ nói chưa, cái tật trùm chăn xem điện thoại suýt nữa thì hại chết con mắt của em đấy!''

Bác sĩ mỉm cười trước màn đối đáp của hai chị em, rồi với tay lấy một lọ thuốc nhỏ mắt có nắp màu đỏ từ trên giá thuốc: ``Tuy nhiên, để biết chính xác mắt trái của cháu cận thật bao nhiêu độ mà không bị hiện tượng co quắp cơ mi đánh lừa, bây giờ bác sĩ sẽ nhỏ thuốc liệt điều tiết. Thuốc này sẽ làm giãn hoàn toàn đồng tử và tạm ngưng chức năng co bóp của thể mi trong vòng hai mươi bốn đến bốn mươi tám tiếng. Cháu sẽ phải ngồi chờ ngoài sảnh khoảng ba mươi phút cho thuốc ngấm hẳn rồi vào đo lại lần cuối.''

Ryuuhou vừa nghe đến đoạn nhỏ thuốc thì mặt lại xanh mét. Con bé bị Kaiho giữ chặt hai vai, ngửa cổ lên trần nhà.

*Tõm! Tõm!*

Hai giọt thuốc trong suốt rơi thẳng vào khóe mắt trái của Ryuuhou. Một cảm giác cay xè, nhồn nhột lan tỏa khiến con bé chớp mắt liên tục, nước mắt sống chảy ròng ròng xuống má, miệng liên tục kêu oai oái: ``Ái ui buốt quá! Bác sĩ ơi sao mắt cháu nó cay xè thế này!''

``Chớp mắt vài cái là hết cay thôi cháu,'' bác sĩ dặn dò thêm một câu cực kỳ quan trọng. ``Lưu ý nhé, khi thuốc ngấm, đồng tử giãn to ra thì cháu sẽ chỉ nhìn được những vật ở xa thôi. Còn tất cả những thứ ở cự ly gần dưới một cánh tay - đặc biệt là sách vở, chữ viết hay màn hình điện thoại - cháu sẽ hoàn toàn không thể nhìn rõ được đâu đấy.''

\ct

Trong lúc dẫn Ryuuhou ra hàng ghế chờ ngoài sảnh lớn để chờ thuốc ngấm, phòng khám lúc này không có thêm bệnh nhân nào mới. Thấy bác sĩ đang rảnh tay, Suzuno bèn thỏ thẻ ngỏ lời: ``Bác sĩ ơi, trong lúc chờ Ryuuhou ngấm thuốc, bác sĩ có thể tiện thể kiểm tra mắt cho mấy anh chị em cháu luôn được không ạ?''

``Được chứ, vào đây cả đi các cháu,'' vị bác sĩ vui vẻ gật đầu.

Thế là cả bốn người chúng tôi lần lượt ngồi vào máy đo.

Bé út Naomi được tôi bế ngồi lên ghế. Con bé háo hức nhìn vào bảng chữ cái vòng tròn hở miệng, bàn tay nhỏ xíu chỉ vanh vách không sai một hình nào. Thị lực của Naomi đạt mười phần mười, đôi mắt trẻ thơ trong veo không tì vết.

Kế đến là Suzuno và Kaiho. Suzuno thường ngày đọc nhiều sách và may vá nhưng giữ khoảng cách rất chuẩn mực, mắt sáng rõ mười phần mười. Còn Kaiho, dù làm việc liên tục với máy tiện, phoi sắt và bản vẽ kỹ thuật cơ khí phức tạp, nhưng nhờ thói quen luôn đeo kính bảo hộ tiêu chuẩn và sinh hoạt điều độ, thị lực hai mắt của cô em song sinh cũng đạt mức tuyệt đối mười phần mười, phản xạ thị giác cực kỳ sắc bén.

Người bước lên ghế cuối cùng là tôi.

Thành thật mà nói, cả Kaiho lẫn Suzuno đều nhìn tôi với ánh mắt đầy hoài nghi. Ai cũng biết tính chất công việc của một kỹ sư phần mềm kiêm quản lý dự án công nghệ như tôi: ban ngày dán mắt vào hai chiếc màn hình máy tính lớn ở văn phòng Tokyo, tối về lại cặm cụi gõ mã nguồn và kiểm tra hệ thống đến tận nửa đêm. Một người ngày đêm tiếp xúc với ánh sáng xanh đặc quánh như tôi, nếu không cận vài ba độ thì chí ít mắt cũng phải mỏi mệt mờ nhòe.

Ngay cả Ryuuhou đang ngồi nhắm nghiền mắt ở ghế chờ cũng hí hửng nói với vào: ``Phen này thể nào anh Kuon cũng cận lòi mắt cho mà xem! Anh cày máy tính còn kinh hơn em xem điện thoại gấp mười lần!''

Tôi mỉm cười không đáp, bình thản tựa cằm vào máy đo khúc xạ rồi nhìn lên bảng kiểm tra thị lực trên tường.

Bác sĩ bắt đầu trượt bảng chữ từ hàng số năm xuống hàng số tám, rồi số chín, và cuối cùng dừng lại ở hàng chữ nhỏ li ti dưới đáy cùng, hàng số mười: ``Đọc cho tôi dòng cuối cùng nào chàng trai.''

Tôi hít một hơi, đọc mạch lạc không vấp một chữ: ``E, P, T, O, L, C, F, D.''

Vị bác sĩ nhướng mày ngạc nhiên, đổi sang bảng chữ ngẫu nhiên khác để thử lại. Tôi vẫn đọc vanh vách từng ký tự một cách dễ dàng.

Bác sĩ tháo máy đo, nhìn màn hình hiển thị thông số rồi bật cười trầm trồ: ``Thật đáng nể đấy! Mắt của cậu thanh niên này đạt mười phần mười ở cả hai mắt. Cậu chỉ bị loạn thị sinh lý cực nhẹ, khoảng không phẩy hai lăm đi-ốp ở trục ngang giác mạc, nhưng đây là cấu trúc bẩm sinh hoàn toàn bình thường, thị lực tự nhiên tự bù đắp được mà không gây mỏi mắt hay giảm tầm nhìn. Nhìn chung, hệ thống điều tiết và võng mạc của cậu cực kỳ khỏe mạnh!''

Tôi mỉm cười gật đầu, nhớ lại chuyện xưa liền tiếp lời: ``Hồi học cấp ba ôn thi tốt nghiệp, cháu cũng từng bị phen hú hồn y như con bé này. Đợt khám sức khỏe định kỳ ở trường năm lớp mười hai, mắt phải của cháu bỗng nhiên bị báo chỉ còn sáu phần mười, làm cả nhà lo sốt vó. Nhưng sau đó đi khám chuyên khoa thì bác sĩ bảo chỉ là 'cận thị giả' do thức đêm cày đề thi quá nhiều. Về sau cháu ngủ đủ giấc, tập thói quen nghỉ ngơi mắt điều độ là thị lực lại hồi phục về mười phần mười như cũ.''

Vị bác sĩ vỗ đùi đánh đét một cái, gật gù tán đồng: ``Đấy, cậu thấy chưa! Chính cậu là nhân chứng sống đấy. Cơ chế của mắt người rất nhạy cảm, biết giữ gìn và cho cơ mắt nghỉ ngơi thì dùng máy tính cả đời cũng chẳng lo thoái hóa.''

``Cái gì cơ?!''

Tiếng kêu thất thanh của Ryuuhou vang dội khắp cả góc hành lang. Con bé ôm lấy một bên mắt, mặt mày méo xệch vì ghen tị: ``Trời đất ơi bất công thế! Sao anh Kuon cày code với chơi game thâu đêm suốt sáng mà mắt vẫn tinh như mắt chim ưng thế kia?! Bác sĩ có nhìn nhầm máy đo không đấy ạ?!''

Kaiho khoanh tay, nhếch mép cười đắc thắng trước sự ấm ức của cô em gái: ``Em đừng có mang cái thói cú vọ phá mắt của mình ra so với anh Kuon! Anh ấy làm việc trên màn hình công nghệ chống lóa cao cấp, ngồi thẳng lưng đúng cự ly chuẩn, phòng làm việc lúc nào cũng bật đủ đèn chiếu sáng và cứ bốn mươi lăm phút lại nghỉ ngơi mắt một lần! Hơn nữa anh ấy từng trải qua bài học sáu phần mười nên biết quý trọng đôi mắt, ai như cái đồ trùm chăn kín mít nằm nghiêng xem phim bộ như em!''

Ryuuhou cứng họng, phụng phịu quay ngoắt mặt đi, không thèm nói thêm lời nào nữa.

\ct

Đúng ba mươi phút sau, thuốc liệt điều tiết đã ngấm hoàn toàn.

Ryuuhou được đưa trở lại phòng đo để kiểm tra khúc xạ lần cuối. Lần này, khi không còn sự can thiệp của các cơn co thắt cơ mi, máy đo đưa ra thông số thực tế chuẩn xác nhất: mắt trái của con bé bị cận thị thực thể ở mức âm không phẩy bảy lăm đi-ốp (-0.75 D), kèm theo độ loạn thị nhẹ không phẩy năm đi-ốp.

``Kết luận cuối cùng nhé,'' bác sĩ tháo găng tay, mỉm cười ôn tồn. ``Độ cận không quá nặng, nhưng vì mắt phải hoàn toàn bình thường nên bắt buộc phải đeo kính điều chỉnh cho mắt trái khi học tập, đọc sách và làm việc trên màn hình. Nếu không đeo, não bộ sẽ tự động bỏ qua tín hiệu từ mắt trái, lâu dần mắt trái sẽ thoái hóa thành nhược thị thật sự đấy.''

Nghe đến việc chắc chắn phải đeo kính, Ryuuhou lại chuẩn bị dẩu môi mè nheo: ``Nhưng mà bác sĩ ơi\dots\ gọng kính ngoài tiệm trông cứ thô kệch thế nào ấy, đeo vào trông già chát chúa\dots''

Suzuno đứng cạnh bỗng khẽ mỉm cười, mở chiếc túi xách vải lấy ra một chiếc hộp nhựa nhỏ đưa cho bác sĩ: ``Thưa bác sĩ, cháu có sẵn một chiếc gọng này ở nhà rồi ạ. Bác sĩ xem giúp cháu liệu có lắp tròng thuốc vào gọng này được không ạ?''

Ryuuhou tròn xoe mắt nhìn chiếc hộp. Đó chính là chiếc kính gọng nhựa retro màu đen nhám viền mỏng cực kỳ cá tính mà con bé từng chắt chiu tiền tiêu vặt mua từ năm ngoái. Hồi đó nó chỉ là một chiếc kính thời trang không độ, được Ryuuhou dùng làm đạo cụ chụp ảnh sống ảo và đeo mỗi khi lên sân khấu biểu diễn cùng ban nhạc rock của trường.

Bác sĩ cầm chiếc gọng lên kiểm tra chốt ốc và độ cong của vành kính, gật đầu tán thưởng: ``Gọng nhựa acetate nguyên khối này rất dẻo dai, ve mũi ôm khít và khoảng cách đồng tử rất chuẩn với khuôn mặt cháu bé. Hoàn toàn lắp tròng thuốc vào được. Mắt trái sẽ lắp tròng điều chỉnh cận loạn, còn mắt phải sẽ lắp tròng không độ có phủ lớp váng lọc ánh sáng xanh để bảo vệ mắt. Vừa thời trang, vừa bảo vệ thị lực tuyệt đối nhé!''

Ryuuhou nhìn chiếc gọng quen thuộc của mình, nét mặt rầu rĩ bỗng tươi tỉnh hẳn lên. Dù sao thì đeo chiếc kính rock-chic quen thuộc của mình vẫn ngầu hơn gấp vạn lần việc phải vác một cái gọng kính kim loại già nua của người lớn.

Sau khi thanh toán chi phí lắp tròng và hẹn chiều muộn quay lại lấy kính, cả năm anh chị em rời khỏi phòng khám để trở về nhà.

\ct

Thế nhưng, cơn ác mộng thực sự của Ryuuhou chỉ mới bắt đầu khi cánh cửa phòng khám khép lại sau lưng.

Vừa bước ra ngoài trời, ánh nắng ban trưa rực rỡ dội xuống khiến con mắt trái đang giãn hết cỡ đồng tử của Ryuuhou chói lòa đến mức không thể mở ra nổi. Con bé vội vàng lấy tay che mắt, nép chặt vào sau lưng tôi như một chú mèo con sợ ánh sáng: ``Chói quá Onii ơi! Mắt trái của em nó cứ lóa trắng xóa như bị ai chiếu đèn pha vào mặt ấy!''

Suzuno vội vàng bung chiếc ô che nắng màu xanh nhạt che chắn cho em gái, trong khi Kaiho nắm chặt lấy cổ tay con bé dắt đi từng bước qua dải công viên.

Và khi về đến căn phòng tầng ba quen thuộc, bi kịch của một ``con nghiện điện thoại'' chính thức bùng nổ.

Tác dụng của thuốc liệt điều tiết khiến thể mi của Ryuuhou hoàn toàn tê liệt. Khi nhìn ra ngoài cửa sổ, dải công viên xanh mướt bên kia bờ sông hay những nhành cây phong đu đưa trong gió hiện lên sắc nét đến từng đường gân lá. Nhưng hễ con bé cúi xuống nhìn bất cứ thứ gì ở cự ly gần dưới một cánh tay, toàn bộ thế giới trước mắt bỗng chốc nhòe nhoẹt, mờ tịt như thể bị phủ một màn sương mù dày đặc.

Cực hình đầu tiên diễn ra ngay trong bữa cơm trưa.

Cầm đôi đũa trên tay, Ryuuhou không thể nào định lượng được khoảng cách của thức ăn trong mâm. Con bé thò đũa gắp một miếng chả lụa nhưng lại gắp trượt vào khoảng không, cắm thẳng đầu đũa vào đĩa dưa chuột muối của Kaiho. Đến khi chan canh, con bé suýt nữa thì đổ cả muôi canh ra ngoài mép bàn nếu không có Suzuno nhanh tay đỡ lấy bát.

Cuối cùng, cô nàng tomboy ngổ ngáo đành phải ngậm ngùi cầm thìa xúc cơm ăn từng miếng một cách vụng về, dưới ánh mắt vừa thương vừa buồn cười của cả nhà.

Nhưng bữa ăn vẫn chưa phải là điều tồi tệ nhất.

Sau bữa trưa, khi bố Ryujin mang chiếc điện thoại thông minh trả lại cho Ryuuhou để con bé kiểm tra tin nhắn học tập từ lớp trưởng, thảm cảnh thực sự mới ập đến.

Ryuuhou háo hức cầm lấy chiếc điện thoại, bấm nút mở khóa màn hình.

Thế nhưng, vừa nhìn vào màn hình ở khoảng cách quen thuộc hai mươi phân, con bé bỗng rú lên một tiếng tuyệt vọng:

``Áaaa! Chữ đâu hết rồi?! Sao toàn là những vệt sáng nhảy múa lung tung thế này?!''

Dưới tác động của đồng tử giãn to, ánh sáng từ màn hình phát ra chói chang như một ngọn đuốc, trong khi các biểu tượng ứng dụng và từng dòng tin nhắn nhòe nhoẹt, biến dạng thành những vệt màu trừu tượng không thể nào đọc nổi. Ryuuhou cố gắng kéo điện thoại ra xa hết cỡ một sải tay, nheo nheo mắt cố đọc thử một dòng thông báo, nhưng chỉ sau chưa đầy mười giây, thái dương bên trái của con bé đã giật thắt từng cơn buốt nhói, nước mắt sống ứa ra tràn mi.

``Thôi ngay đi!'' Kaiho từ bàn bên bước sang, giật phắt chiếc điện thoại cất vào ngăn kéo tủ gỗ rồi khóa lại đánh *CẠCH* một cái. ``Bác sĩ đã dặn là thuốc liệt điều tiết làm mất khả năng nhìn gần, càng cố nhìn càng làm rách cơ mắt đấy! Hai ngày cuối tuần này, cấm tiệt ti vi, điện thoại, máy tính!''

Thế là suốt hai ngày thứ Bảy và Chủ nhật, Ryuuhou phải trải qua một cuộc ``cai nghiện công nghệ'' đau đớn và vật vã chưa từng có trong đời.

Căn phòng tầng ba biến thành một ``nhà tù'' đối với cô nhóc thừa năng lượng. Không thể lướt mạng xã hội, không thể cày phim hoạt hình, không thể xem các video hướng dẫn bấm hợp âm guitar rock, con bé nằm vật vã trên chiếc ghế đệm dài, lăn qua lộn lại từ đầu này sang đầu kia, miệng không ngừng rên rỉ như một người bị giam lỏng:

``Chán quá đi mất thôi\dots\ Ai cứu tôi với\dots\ Tôi sắp mọc rêu trên người rồi đây này\dots''

Bé út Naomi ngồi dưới sàn vẽ tranh, thỉnh thoảng lại ngẩng lên giơ một bức tranh màu sáp nguệch ngoạc lên hỏi: ``Chị Ryuuhou xem Naomi vẽ con khủng long có đẹp không?''

Ryuuhou nheo mắt nhìn đám màu sắc nhòe nhoẹt trước mặt, thở dài thườn thượt: ``Đẹp\dots\ đẹp lắm em ạ, trông giống hệt một củ khoai tây bị nướng cháy\dots''

Đến chiều Chủ nhật, quá bức bối vì không có gì giải trí, Ryuuhou đành phải hạ mình năn nỉ Suzuno ngồi bên cạnh đọc to từng mẩu truyện ngắn trinh thám cho mình nghe. Đến tối, con bé lại lò dò mò lên tầng năm, ngồi bệt xuống cạnh bàn làm việc của tôi, ngoan ngoãn như một chú cún con mè nheo nhờ tôi kể chuyện về những ngày tháng tôi đi học và làm dự án bên Tokyo cho đỡ buồn ngủ.

Suốt bốn mươi tám tiếng đồng hồ không có ánh sáng xanh của màn hình điện thoại, Ryuuhou buộc phải để đôi mắt của mình nghỉ ngơi hoàn toàn. Mỗi buổi chiều, Kaiho lại nghiêm khắc xách tai con bé ra ban công tầng ba, bắt đứng tựa vào lan can đúng ba mươi phút, phóng tầm mắt qua dải công viên ven sông để ngắm nhìn những rặng cây xanh rì và bầu trời hoàng hôn bao la.

Đó có lẽ là hai ngày dài nhất trong cuộc đời của Ryuuhou. Nhưng cũng chính nhờ hai ngày sống chậm bất đắc dĩ ấy, đôi mắt mỏi mệt suốt bao tháng ngày của cô nhóc cuối cùng cũng được hít thở một bầu không khí trong lành, chuẩn bị cho ngày đón nhận một góc nhìn hoàn toàn mới về thế giới xung quanh qua tấm kính nhỏ.

\pov{Hikawa Kaiho}

Chiều thứ Hai, đợt tác dụng của liều thuốc liệt điều tiết rốt cuộc cũng chấm dứt hoàn toàn.

Đồng tử ở mắt trái của Ryuuhou đã co nhỏ trở lại kích thước bình thường, không còn nhạy cảm với ánh sáng ban ngày nữa. Con bé đã có thể nhìn gần trở lại, và chiếc điện thoại thông minh cũng được bố trao trả - tất nhiên là dưới sự giám sát giờ giấc vô cùng ngặt nghèo của cả nhà.

Chiều muộn hôm trước, tôi và anh Kuon đã ghé qua tiệm kính thuốc để lấy chiếc gọng kính nhựa retro lắp tròng mới về. Khi ngồi vào bàn học ở tầng ba làm bài tập, Ryuuhou ngoan ngoãn đeo kính lên sống mũi. Đúng như chẩn đoán của bác sĩ, nhờ có tròng kính điều chỉnh độ cận loạn bên mắt trái, con bé đọc sách và chép bài trơn tru hơn hẳn, không còn bị nhức thái dương hay mờ nhòe chữ viết như trước nữa.

Thế nhưng, đứng từ góc bếp quan sát, tôi nhận ra nét mặt con bé vẫn phảng phất một nỗi muộn phiền giấu kín. Mỗi lần tháo chiếc kính ra đặt lên mặt bàn gỗ, con bé lại lén thở dài một hơi đầy tiếc nuối.

Tôi thừa biết tỏng cô em gái ngổ ngáo của mình đang trăn trở điều gì.

Ryuuhou là đứa em gái mê rock đến cuồng nhiệt. Con bé là cây guitar điện chủ lực của ban nhạc học sinh ở trường, và vào cuối tháng này, ban nhạc của nó sẽ có một buổi biểu diễn trực tiếp trên sân khấu lớn trong lễ hội giao lưu mùa xuân.

Đối với một tay guitar chơi nhạc rock đầy lửa, việc phải vác một cặp kính gọng lên mặt biểu diễn chẳng khác nào một thảm họa thực sự. Khi bước lên sân khấu, đắm mình vào những đoạn solo cuồng nhiệt, vừa gảy đàn vừa nhảy nhót, headbang theo từng nhịp trống dồn dập, mồ hôi túa ra như tắm sẽ khiến chiếc gọng kính trơn tuột trên sống mũi. Chỉ cần sơ sẩy một chút, kính có thể văng ra sàn sân khấu vỡ tan tành; mà giá cho mỗi tròng cũng chẳng rẻ gì cho cam. Chưa kể, vành gọng kính sẽ thu hẹp tầm nhìn ngoại vi khi con bé muốn quan sát các thành viên khác trong ban nhạc. Và quan trọng nhất đối với một cô bé mười sáu tuổi: chiếc gọng kính sẽ phá hỏng hoàn toàn phong cách bụi bặm, phóng khoáng và cool ngầu mà con bé dày công xây dựng bấy lâu nay.

Là chị cả trong nhà, tôi có thể nghiêm khắc, có thể đập bàn quát mắng khi em mình làm điều sai trái, nhưng tôi tuyệt đối không bao giờ để em gái mình phải chịu ấm ức hay mất đi sự tự tin trên sân khấu ước mơ của nó.

Đó chính là lý do vì sao vào sáng thứ Bảy tuần trước, sau khi anh Kuon dắt tụi nhỏ rời khỏi phòng khám, tôi đã cố tình nán lại phía sau để gặp riêng vị bác sĩ nhãn khoa.

Bằng số tiền tiêu vặt tích cóp từ khoản phụ cấp làm thêm ở xưởng cơ khí của gia đình, tôi đã nhờ bác sĩ tư vấn và đặt riêng cho Ryuuhou một bộ kính áp tròng mềm chuyên dụng cao cấp. Đó là loại kính áp tròng dùng theo tháng với độ ngậm nước cao, mỏng nhẹ và thoáng khí tối đa, được đo cắt chuẩn xác theo thông số cận loạn của mắt trái con bé, đi kèm một bộ khay đựng và chai nước ngâm sinh lý sát khuẩn.

Tôi bước lại gần chiếc bàn trà ở tầng ba, ném nhẹ chiếc hộp nhỏ bọc giấy gói cẩn thận xuống trước mặt Ryuuhou: ``Này, cầm lấy.''

Ryuuhou đang ủ rũ gục đầu trên bàn, giật mình ngẩng lên ngơ ngác: ``Gì đây chị Kaiho?''

``Mở ra thì biết.''

Con bé tò mò bóc lớp giấy bọc bên ngoài. Khi nhìn thấy chiếc hộp đựng kính áp tròng nhỏ gọn cùng lọ dung dịch ngâm kính chuyên dụng, đôi mắt một mí sắc sảo của Ryuuhou bỗng chốc tròn xoe, miệng há hốc kinh ngạc: ``Kính\dots\ kính áp tròng mềm sao?! Chị mua cho em từ bao giờ thế này?!''

Tôi khoanh hai tay trước ngực, cố tình giữ nét mặt lạnh tanh: ``Chị đoán trước rồi. Không chuẩn bị cái này thì tuần tới lên sân khấu lễ hội trường, em tính vừa gảy guitar vừa lấy ngón tay đẩy gọng kính à? Hay tính nhảy cho đã đời rồi để kính văng xuống sàn, mù dở ngã nhào xuống đầu khán giả?''

Nghe những lời cộc cằn nhưng đong đầy sự thấu hiểu ấy, Ryuuhou sững sờ mất vài giây. Rồi bất thình lình, con bé nhảy cẫng lên khỏi ghế, hai mắt long lanh như muốn lao tới ôm chầm lấy tôi: ``Ôi chị Kaiho! Chị là người chị tuyệt vời nhất trần đời! Em yêu chị chết mất thôi!''

``Tránh xa ra, bớt nịnh nọt đi!'' tôi vội lùi lại một bước, đưa tay chặn trán con bé, hắng giọng nhắc nhở. ``Mua cho là một chuyện, đeo được vào mắt hay không lại là chuyện khác đấy. Mau vào nhà vệ sinh rửa tay thật sạch rồi ra đây chị hướng dẫn.''

Thế là căn phòng tầng ba lại được một phen náo nhiệt với màn tập đeo kính áp tròng lần đầu tiên trong đời của Ryuuhou.

Đó quả thực là một màn vật lộn vừa hài hước vừa dở khóc dở cười.

Ryuuhou ngồi ngay ngắn trước chiếc gương tròn nhỏ đặt trên bàn trà. Con bé dùng chiếc nhíp nhỏ gắp miếng thấu kính mỏng manh, trong suốt như một giọt nước đặt lên đầu ngón tay trỏ. Thế nhưng, cứ mỗi lần con bé đưa ngón tay lại gần sát mắt, phản xạ tự nhiên của cơ thể lại khiến mí mắt nó nhắm tịt lại như kéo rèm, thấu kính mềm dính ngược vào lông mi rồi rơi tuột xuống mu bàn tay.

``Ái ui sợ quá!'' Ryuuhou co rúm người lại, nhăn tít mặt mũi kêu than. ``Ngón tay chọc thẳng vào lòng đen thế này trông ghê chết đi được! Cứ như tự chọc mù mắt mình ấy!''

Suzuno đứng bên cạnh nhẹ nhàng cầm chiếc khăn sạch, kiên nhẫn cổ vũ: ``Bình tĩnh nào Ryuuhou. Hít một hơi thật sâu, dùng ngón tay trái giữ chặt mí mắt trên sát chân mày, mắt nhìn thẳng vào gương đừng đảo qua đảo lại\dots''

Bé út Naomi thì ghé sát cằm lên mép bàn, hai mắt chớp chớp tò mò theo dõi từng cử động như thể đang xem một nhà ảo thuật chuẩn bị làm xiếc.

Thấy con bé loay hoay đến lần thứ tư mà vẫn chưa đeo nổi, mồ hôi lấm tấm trên trán, tôi sốt ruột không nhịn được nữa: ``Tránh ra để chị làm cho! Có con mắt thôi mà làm như sắp ra trận đánh nhau không bằng!''

Tôi rửa tay sát khuẩn lại một lần nữa bằng xà phòng y tế rồi bước tới trước mặt Ryuuhou. Tôi bảo con bé ngửa cổ ra sau, mắt nhìn thẳng vào tôi. Bằng đôi bàn tay quen với việc căn chỉnh từng chi tiết cơ khí tinh vi dưới một phần mười milimét, tôi nhẹ nhàng dùng hai ngón tay cố định chắc chắn mí trên và mí dưới của con bé, rồi khéo léo đặt thấu kính mềm mại áp nhẹ vào bề mặt nhãn cầu mắt trái.

*NHẤP*

Thấu kính êm ái ôm khít lấy giác mạc.

``Được rồi đấy, chớp mắt nhẹ vài cái xem nào,'' tôi buông tay ra, lùi lại một bước.

Ryuuhou chớp chớp mắt hai ba cái. Cảm giác man mát, dịu nhẹ lan tỏa khắp nhãn cầu. Con bé đứng bật dậy, chớp mắt lia lịa, vẻ mặt đầy bỡ ngỡ: ``Ơ\dots\ không hề thấy đau hay cộm gì luôn! Cứ như không có gì trong mắt ấy!''

``Bước ra cửa sổ xem thử đi,'' tôi khẽ hất cằm về phía khung kính lớn nhìn ra bờ sông.

Ryuuhou háo hức chạy vội ra bên khung cửa sổ tầng ba, nơi ánh hoàng hôn rực rỡ đang nhuộm vàng cả một góc trời Anwa.

Cô nhóc đứng sững lại giữa làn gió xuân nhè nhẹ thổi qua khe cửa sổ.

Đôi mắt của con bé mở to, phản chiếu ánh chiều tà lấp lánh. Không còn chiếc gọng kính vướng víu nào cản trở tầm nhìn ngoại vi, không còn cảm giác nặng nề đè lên sống mũi hay vành tai, và quan trọng nhất, toàn bộ không gian trước mắt mở ra với một chiều sâu ba chiều sống động đến nghẹt thở.

Từng rặng cây phong đu đưa bên dải công viên bờ sông, từng chiếc lá xanh non dập dờn trong gió, dòng nước sông trước cửa nhà lấp lánh ánh bạc, con dốc thoải trước cửa nhà và cả những mái ngói rêu phong của khu phố cổ xa xa\dots\ tất cả đều hiện lên rõ ràng, sắc nét đến từng đường nét chi tiết mà trước đây mắt trái mờ tịt của con bé chưa từng được cảm nhận trọn vẹn.

Ryuuhou đứng lặng người rất lâu, dường như bị choáng ngợp trước vẻ đẹp chân thực của thế giới qua tấm kính nhỏ vô hình ấy.

Đúng lúc đó, tiếng bước chân quen thuộc vang lên trên cầu thang. Anh Kuon vừa tan ca làm trực tuyến, bước xuống tầng ba với chiếc ấm trà trên tay. Nhìn thấy Ryuuhou đang đứng ngây người bên cửa sổ, anh nhướng mày ngạc nhiên: ``Đeo được kính áp tròng rồi à Ryuu?''

Ryuuhou không quay đầu lại ngay. Con bé vẫn dán mắt vào khoảng không gian bao la ngoài kia, rồi bỗng cất tiếng hỏi khẽ: ``Anh Kuon này.''

Anh Kuon dừng chân bên bàn trà: ``Hửm?''

``Cái cây phong ở góc công viên ven sông đằng kia\dots\ có bao nhiêu cành lớn vậy anh?''

Anh Kuon hơi bất ngờ trước câu hỏi kỳ lạ của em gái. Anh bước lại gần khung cửa sổ, nheo mắt nhìn theo hướng tay chỉ của con bé, lẩm nhẩm đếm một lúc rồi trả lời: ``Khoảng\dots\ mười bảy cành lớn? Sao tự nhiên lại hỏi thế?''

Lúc này, Ryuuhou mới từ từ quay người lại.

Trên gương mặt của cô nàng tomboy ngổ ngáo thường ngày không còn nét nhăn nhó hay vẻ bướng bỉnh quen thuộc, mà thay vào đó là một nụ cười -  rất khẽ, rất sáng, và rạng rỡ từ tận đáy lòng:

``Không có gì đâu ạ. Em chỉ muốn thử xem mình đếm có đúng không thôi.''

Tôi đứng tựa lưng vào mép cửa ra vào, khoanh tay nhìn nụ cười trong trẻo của đứa em gái nhỏ. Bao nhiêu bực bội, bao nhiêu lo lắng và cả những lời mắng mỏ gay gắt suốt mấy ngày qua bỗng chốc tan biến như mây khói theo làn gió xuân. Khóe môi tôi bất giác cũng khẽ cong lên một nụ cười nhẹ nhõm.

\omk

Chiều ngày hôm sau, dưới xưởng cơ khí của gia đình chúng tôi tại ngoại ô.

Sau giờ tan học, Ryuuhou không còn nằm dài trùm chăn lướt điện thoại như trước nữa. Con bé đeo chiếc kính áp tròng, ngoan ngoãn mặc chiếc áo bảo hộ lao động mỏng rồi tung tăng xuống xưởng phụ giúp tôi và bố Ryujin quét dọn mạt sắt và lau dầu bảo dưỡng cho dàn máy tiện.

Đang cặm cụi dùng chổi lông quét mạt phôi trên bàn gá kim loại, con bé bỗng dừng phắt lại. Ryuuhou đứng chôn chân giữa xưởng, ngửa cổ nhìn chằm chằm lên góc ô cửa sổ kính mờ đóng bụi ở phía góc tường cao.

Tôi đang cầm thước cặp đo đường kính trục thép, thấy con bé đứng bất động như tượng suốt năm phút đồng hồ liền thắc mắc cất tiếng: ``Này Ryuu! Em đứng ngây ra đấy làm cái trò gì vậy? Không mau quét nốt đống mạt sắt đi cho bố tắt máy.''

``Chị Kaiho ra đây mà xem này!'' Ryuuhou vẫy vẫy tay, giọng đầy vẻ thích thú. ``Em đang nhìn con nhện!''

Tôi nhíu mày bước lại gần: ``Con nhện?''

``Đúng rồi! Ở góc trên cùng của khung cửa sổ có một con nhện bé tí xíu đang giăng tơ kìa!'' Ryuuhou chỉ tay lên góc tường, mắt lấp la lấp lánh như vừa phát hiện ra kho báu. ``Bình thường đứng ở khoảng cách xa thế này, mắt trái em nhòe nhoẹt chả thấy gì ngoài một đốm đen mờ tịt. Thế mà hôm nay em nhìn thấy rõ từng cái chân của nó, thấy cả sợi tơ mảnh li ti óng ánh dưới ánh đèn tuýp nữa! Trông hay cực kỳ luôn!''

Tôi ngước nhìn theo ngón tay con bé. Quả thực ở góc khung kính mờ có một chú nhện nhỏ xíu đang cần mẫn dệt lưới.

Nhìn vẻ mặt phấn khích như một đứa trẻ con của cô em gái mười bảy tuổi, tôi chỉ biết lắc đầu phì cười. Tôi tiện tay rút một chiếc giẻ lau máy sạch sẽ dúi thẳng vào tay nó: ``Rõ mắt rồi thì lau nốt cái bàn gá máy phay kia đi! Ở đấy mà ngắm nhện với giăng tơ, khéo mạt sắt nó bay vào kính áp tròng bây giờ!''

``Tuân lệnh đội trưởng Kaiho!'' Ryuuhou cười tít mắt, cầm lấy chiếc giẻ lau rồi thoăn thoắt bắt tay vào việc, tràn đầy năng lượng của một tuổi trẻ sáng ngời.

\end{document}